\documentclass[tikz,border=8pt]{standalone}
\usepackage{tikz}
\usepackage{amsmath}
\usepackage{amssymb}
\usepackage{pgfplots}
\pgfplotsset{compat=1.18}
\usepackage{ctex}
\usetikzlibrary{calc,positioning,arrows.meta,trees}

% LC 199/102/103 BFS 三种变体对比
% 树：[3,9,20,null,null,15,7]（5节点）
% 左：199 右视图 [3,20,7]
% 中：102 层序 [[3],[9,20],[15,7]]
% 右：103 锯齿 [[3],[20,9],[15,7]]

\begin{document}
\begin{tikzpicture}[
  every node/.style={},
  level distance=1.3cm,
  level 1/.style={sibling distance=2.4cm},
  level 2/.style={sibling distance=1.4cm},
  edge from parent/.style={draw,->,>=Stealth}
]

%% ── 标题 ──────────────────────────────────────────────────
\node[font=\large\bfseries] at (8.5, 1.6) {BFS 三种变体对比（同一棵树）};
\node[font=\small,gray] at (8.5, 1.1) {树 \texttt{[3,9,20,null,null,15,7]}};

%% ══════════════════════════════════════════════════
%% 左列：LC 199 右视图
%% ══════════════════════════════════════════════════
\node[font=\small\bfseries,blue!70!black] at (2.5, 0.5) {LC 199 右视图};

\node[draw,circle,minimum size=0.72cm,fill=blue!10,font=\bfseries] (l-root) at (2.5,0) {3}
  child {
    node[draw,circle,minimum size=0.72cm,fill=blue!10,font=\bfseries] {9}
  }
  child {
    node[draw,circle,minimum size=0.72cm,fill=blue!10,font=\bfseries] (l-n20) {20}
    child {
      node[draw,circle,minimum size=0.72cm,fill=blue!10,font=\bfseries] {15}
    }
    child {
      node[draw,circle,minimum size=0.72cm,fill=blue!10,font=\bfseries] (l-n7) {7}
    }
  };

% 右视图高亮：3, 20, 7 用橙色圆圈
\draw[orange!80!black,very thick,rounded corners=2pt] (l-root) circle (0.42cm);
\draw[orange!80!black,very thick,rounded corners=2pt] (l-n20) circle (0.42cm);
\draw[orange!80!black,very thick,rounded corners=2pt] (l-n7) circle (0.42cm);

% 右视图结果框
\node[draw=gray!50,fill=orange!5,rounded corners=3pt,inner sep=5pt,
      font=\small,align=center]
  at (2.5,-4.0)
  {\textbf{右视图}\\每层最右节点\\$[3,\ 20,\ 7]$};

%% ══════════════════════════════════════════════════
%% 中列：LC 102 层序
%% ══════════════════════════════════════════════════
\node[font=\small\bfseries,green!60!black] at (8.5, 0.5) {LC 102 层序遍历};

\node[draw,circle,minimum size=0.72cm,fill=blue!10,font=\bfseries] (m-root) at (8.5,0) {3}
  child {
    node[draw,circle,minimum size=0.72cm,fill=blue!10,font=\bfseries] (m-n9) {9}
  }
  child {
    node[draw,circle,minimum size=0.72cm,fill=blue!10,font=\bfseries] (m-n20) {20}
    child {
      node[draw,circle,minimum size=0.72cm,fill=blue!10,font=\bfseries] (m-n15) {15}
    }
    child {
      node[draw,circle,minimum size=0.72cm,fill=blue!10,font=\bfseries] (m-n7) {7}
    }
  };

% 层序：逐层用横线分隔标注
\draw[green!50!black,dashed,thin] (6.5,-0.65) -- (10.5,-0.65) node[right,font=\tiny,green!50!black] {L0};
\draw[green!50!black,dashed,thin] (6.5,-1.95) -- (10.5,-1.95) node[right,font=\tiny,green!50!black] {L1};
\draw[green!50!black,dashed,thin] (6.5,-3.25) -- (10.5,-3.25) node[right,font=\tiny,green!50!black] {L2};

\node[draw=gray!50,fill=green!5,rounded corners=3pt,inner sep=5pt,
      font=\small,align=center]
  at (8.5,-4.0)
  {\textbf{层序结果}\\$[[3],\ [9,20],\ [15,7]]$};

%% ══════════════════════════════════════════════════
%% 右列：LC 103 锯齿形
%% ══════════════════════════════════════════════════
\node[font=\small\bfseries,purple!70!black] at (14.5, 0.5) {LC 103 锯齿遍历};

\node[draw,circle,minimum size=0.72cm,fill=blue!10,font=\bfseries] (r-root) at (14.5,0) {3}
  child {
    node[draw,circle,minimum size=0.72cm,fill=blue!10,font=\bfseries] (r-n9) {9}
  }
  child {
    node[draw,circle,minimum size=0.72cm,fill=blue!10,font=\bfseries] (r-n20) {20}
    child {
      node[draw,circle,minimum size=0.72cm,fill=blue!10,font=\bfseries] (r-n15) {15}
    }
    child {
      node[draw,circle,minimum size=0.72cm,fill=blue!10,font=\bfseries] (r-n7) {7}
    }
  };

% 锯齿：Level1 从右到左箭头，Level2 从左到右
\draw[->,>=Stealth,purple!60,very thick] (r-n20) to[bend right=20] (r-n9);
\node[font=\tiny,purple!70!black] at (13.5,-1.4) {←右到左};
\draw[->,>=Stealth,purple!40,thick] (r-n15) to[bend right=10] (r-n7);
\node[font=\tiny,purple!70!black] at (14.5,-3.0) {左到右→};

\node[draw=gray!50,fill=purple!5,rounded corners=3pt,inner sep=5pt,
      font=\small,align=center]
  at (14.5,-4.0)
  {\textbf{锯齿结果}\\奇数层反向\\$[[3],\ [20,9],\ [15,7]]$};

%% ── 分隔线 ───────────────────────────────────────────────
\draw[gray!40,very thin,dashed] (5.5,-4.8) -- (5.5,1.8);
\draw[gray!40,very thin,dashed] (11.5,-4.8) -- (11.5,1.8);

%% ── 对比说明 ─────────────────────────────────────────────
\node[draw=gray!50,fill=white,thin,rounded corners=3pt,
      font=\scriptsize,inner sep=5pt,align=left]
  at (8.5,-5.5)
  {三者区别：右视图取每层末尾元素；层序按层收集；锯齿在 Level 1,3,\ldots 反转列表};

\end{tikzpicture}
\end{document}
